% =====================================================================
\section{Method}
\label{sec:method}
% =====================================================================
% DRAFT: follows plan/method_plan.md (revised 2026-10-06). Architecture sizes, losses and schedules are
% filled in once v0/v1 train. Keep in sync with the code.

\begin{figure*}[t]
  \centering
  \resizebox{\textwidth}{!}{% Method overview (Sec. 3). TikZ on an explicit grid; frozen = grey, trained = blue, data = white.
\begin{tikzpicture}[
  font=\scriptsize, >=Latex,
  data/.style={draw, rounded corners=1pt, fill=white, align=center, inner sep=2pt, minimum height=7mm},
  frozen/.style={draw, rounded corners=2pt, fill=gray!15, align=center, inner sep=2pt, minimum height=7mm},
  train/.style={draw=blue!60!black, rounded corners=2pt, fill=blue!8, align=center, inner sep=2pt, minimum height=8mm},
  op/.style={draw, circle, inner sep=1pt, fill=white},
  lab/.style={font=\scriptsize\itshape, text=black!70, fill=white, inner sep=1pt},
  ar/.style={->, thick, black!75},
]
% ------------------------------------------------------------------ per image
\node[font=\scriptsize\bfseries, anchor=west] at (-1.0, 1.0) {per image};
\node[data, minimum width=18mm] (x) at (0, 0) {native input $x$\\[-1pt]{\tiny $6720{\times}4480$}};
\node[frozen] (D) at (2.7, 0) {deblurrer $\mathcal{D}$\\[-1pt]{\tiny any method, at $\times4$}};
\node[data] (a) at (5.3, 0) {anchor $a$\\[-1pt]{\tiny $1680{\times}1120$}};
\node[frozen] (dino) at (7.8, 0) {DINOv2\\[-1pt]{\tiny tokens of $a$}};
\node[data, minimum width=44mm] (mem) at (12.3, 0) {in-focus exemplar memory $\mathcal{M}$\\[-1pt]
  {\tiny keys: anchor tokens in focus \;$\cdot$\; values: native crops of $x$}};
\node[data, minimum width=18mm] (d) at (0, -1.3) {DP map $d$\\[-1pt]{\tiny $|$disparity$|$, confidence}};
\draw[ar] (x) -- node[above, lab] {$\downarrow_4$} (D);
\draw[ar] (D) -- (a);
\draw[ar] (a) -- (dino);
\draw[ar] (dino) -- node[above, lab] {keys} (mem);
\draw[ar] (x.north) -- ++(0, 0.35) -| node[pos=0.25, above, lab] {native values} ([xshift=-12mm]mem.north);
\draw[ar] (d.east) -- node[pos=0.5, above, lab] {in-focus regions select the keys} ([xshift=-12mm]mem.south |- d.east)
  -- ([xshift=-12mm]mem.south);
% ------------------------------------------------------------------ per tile
\node[font=\scriptsize\bfseries, anchor=west] at (-1.0, -2.35) {per $512{\times}512$ tile $T$};
\node[data, minimum width=18mm] (xt) at (0, -3.3) {$x_T$, $d_T$, $a_T$\\[-1pt]{\tiny crops of $x$, $d$, $a$}};
\draw[ar] (d.south) -- (xt.north);
\node[train] (cond) at (2.5, -3.3) {conditioning\\trunk};
\node[train, minimum width=36mm, minimum height=13mm] (bb) at (5.9, -3.3) {};
\node[anchor=north, align=center] at (bb.north) {$\times4$ backbone on $a_T$ {\tiny (HAT-L init.)}};
\foreach \i in {1,...,6} {\node[draw=blue!60!black, fill=white, minimum width=3.6mm, minimum height=3.2mm, inner sep=0pt]
  (s\i) at ([xshift={(\i-3.5)*5.4mm}, yshift=-1.0mm]bb.center) {};}
\node[font=\tiny, anchor=south] at (bb.south) {SFT after every stage};
\node[train] (xa) at (9.5, -3.3) {exemplar\\cross-attention\\[-1pt]{\tiny top-$K$ + null}};
\node[train] (ref) at (11.9, -3.3) {native\\refinement};
\node[op] (lock) at (13.55, -3.3) {$\Pi_a$};
\node[data] (yt) at (14.55, -3.3) {$\hat y_T$};
\node[data] (y) at (16.0, -3.3) {$\hat y$\\[-1pt]{\tiny blend, then $\Pi_a$}};
\draw[ar] (xt) -- (cond);
\draw[ar] (cond) -- node[above=1mm, lab] {add, SFT} (bb);
\draw[ar] (bb) -- (xa);
\draw[ar] (xa) -- (ref);
\draw[ar] (ref) -- node[above, lab] {$\Delta_T$} (lock);
\draw[ar] (lock) -- (yt);
\draw[ar] (yt) -- (y);
\draw[ar, blue!60!black] ([xshift=12mm]mem.south) |- ([yshift=4mm]xa.north) node[pos=0.75, above, lab] {exemplars retrieved for the tile's tokens} -- (xa.north);
\draw[ar] ([xshift=-3mm]xt.south) -- ++(0, -0.55) -| node[pos=0.25, below, lab] {$x_T$, $d_T$: copy / deconvolve} (ref.south);
\draw[ar] ([xshift=3mm]xt.south) -- ++(0, -0.95) -| node[pos=0.25, below, lab] {$a_T$} (lock.south);
% ------------------------------------------------------------------ legend
\node[frozen, minimum height=3mm, minimum width=6mm, inner sep=1pt, font=\tiny, anchor=west] (l1) at (-1.0, -5.0) {frozen};
\node[train, minimum height=3mm, minimum width=6mm, inner sep=1pt, font=\tiny, right=2mm of l1] (l2) {trained};
\node[font=\tiny, right=4mm of l2, anchor=west] {$\Pi_a$: anchor lock $\Delta + U(a - \Delta{\downarrow_4})$ with
  $(e{\uparrow}){\downarrow_4}=e$, so that $\hat y{\downarrow_4} = a$};
\end{tikzpicture}
}
  \caption{\textbf{Overview.} \emph{Per image:} any deblurring method $\mathcal{D}$ produces the anchor $a$ at
  $4\times$ lower resolution; DINOv2 features of the anchor index the in-focus regions (from the dual-pixel map $d$)
  into a memory whose values are native crops of the observation $x$. \emph{Per tile:} the native crop $x_T$ and
  $d_T$ condition every stage of a $\times4$ network applied to the anchor tile; exemplars retrieved for the tile's
  tokens enter through cross-attention (with a null option); a native refinement sees $x_T$ directly. The anchor lock
  $\Pi_a$ keeps the result's $\times4$ band equal to the anchor, per tile and again after blending.}
  \label{fig:method}
\end{figure*}

Given a native-resolution defocused image $x \in \mathbb{R}^{H\times W\times 3}$, an existing deblurring
method $\mathcal{D}$ is applied to its $4\times$ downscaled version and yields the \emph{anchor}
$a = \mathcal{D}(x{\downarrow_4})$, which has the structure of the all-in-focus image but no native
detail. \method produces the native all-in-focus estimate $\hat{y}$ from $x$, $a$ and a blur map $d$
giving the local amount of defocus (\cref{fig:method}). It operates on $512\times512$ tiles $T$ with overlap and blending,
but every tile has access to a memory $\mathcal{M}$ of in-focus exemplars built once per image
(\cref{sec:memory}).

\subsection{Anchor-locked upsampling}
\label{sec:lock}
The baselines show that the anchor's structure is worth keeping exactly: bicubic upsampling preserves
it (the result, downscaled again, scores like the anchor itself), whereas generative upsamplers alter
it and lose fidelity even at the anchor's resolution (\cref{sec:exp_results}). We therefore let the
network contribute only what the anchor cannot express. With $G$ the network, $a{\uparrow}$ the
bicubically upsampled anchor, ${\downarrow_4}$ the area downscaling used throughout and ${\uparrow}$ an
upsampling with $(e{\uparrow}){\downarrow_4} = e$ (bicubic, corrected by its own nearest-neighbour residual),
\begin{equation}
  \begin{aligned}
    \Delta_T &= G\big(x_T,\, a{\uparrow}_T,\, d_T;\, \mathcal{M}\big),\\
    \hat{y}_T &= a{\uparrow}_T + \Delta_T - (\Delta_T{\downarrow_4}){\uparrow},
  \end{aligned}
  \label{eq:output}
\end{equation}
so that $\hat{y}{\downarrow_4} = a$ exactly, as in one step of iterative
back-projection~\cite{irani1991backprojection}; after tiles are blended, the same projection is applied once more
to the whole image. The output is as faithful as the anchor at the
benchmark resolution by construction, any improvement of the anchor carries over directly, and the
network is free to spend its capacity on the native high band. The price is that structural errors of
the anchor are inherited; we leave correcting them to the low-resolution method.

\subsection{Where native detail comes from}
\label{sec:regimes}
Defocus varies across the image, and with it the source of the missing high band. \emph{In focus}, the
observation already holds it: the network should reproduce $x$, lightly denoised (the f/22 targets
are less noisy than the f/4 inputs). Under \emph{mild and moderate defocus}, the circle of confusion
attenuates but does not remove the frequencies above the anchor's band (the first zero of a disc
blur of radius $\rho$ native pixels lies at $0.61/\rho$ cycles per pixel, above the anchor's Nyquist
frequency of $1/8$ for $\rho \lesssim 5$): the detail can be recovered from $x$ itself by local
deconvolution, a regime that neither native deblurring nor the low-resolution pipelines exploit.
Under \emph{strong defocus} it is gone, and must come from elsewhere: from in-focus exemplars of the
same image when they exist (\cref{sec:memory}), and from the network's prior otherwise. $G$ receives
the blur map so that it can tell these regimes apart. We use the dual-pixel disparity directly as $d$,
which is monotonic in the defocus and needs no calibration, only lightly smoothed so that depth edges stay sharp,
together with its matching confidence, which separates ``in focus'' from ``no texture to measure''; converting it to a circle-of-confusion
radius or to per-band availability maps is evaluated as an ablation. When dual-pixel data are
unavailable, $d$ is predicted from the agreement between the anchor and the downscaled
observation~\cite{zhuo2011defocusmap}. \todo{predicted-map head and its training.}

\subsection{In-focus exemplar memory}
\label{sec:memory}
Inside a strongly defocused tile neither the observation nor the anchor carries native texture, but
the same material is often in focus elsewhere in the photograph~\cite{glasner2009single,zontak2011internal}.
\emph{Features.} Defocused regions and their in-focus counterparts cannot be matched at native
resolution, where one is blurred and the other sharp. In the anchor they are comparable, since the
anchor is deblurred everywhere. We compute DINOv2~\cite{oquab2024dinov2} features once on the whole
anchor; at the anchor resolution each token covers about $56\times56$ native pixels.
\emph{Memory.} Keys are the anchor tokens of in-focus regions (small $|d|$), excluding flat ones; the
corresponding values are the native observation crops under them, encoded at three scales to account
for perspective.
\emph{Retrieval.} The tokens of the defocused part of a tile are the queries. For each, the $K$ most
similar keys by cosine similarity, their similarity scores and a learned null entry are passed to $G$,
which attends to them,
\begin{equation}
  r = \operatorname{softmax}\!\big(q^\top [k_{j_1},\dots,k_{j_K}, k_\varnothing] / \tau\big),
  \label{eq:relevance}
\end{equation}
and receives the transferred features $\sum_i r_i v_{j_i}$. Retrieval therefore needs recall rather
than precision; relevance is learned. The null entry lets $G$ abstain and fall back to its prior where
nothing relevant is in focus, and $\max_i r_i$ is a relevance map that we use for analysis. Because
the memory is global, every tile draws on the whole image, and the result does not depend on the tile
size. During training the memory is built from the whole image while the loss is computed on a crop,
and exemplars overlapping the crop are dropped so that the model learns to transfer rather than copy
in place.

\subsection{Network and training}
\label{sec:variants}
\paragraph{Conditioning.} The native input, rearranged onto the anchor grid, and $d$ are encoded by a small
trunk whose features modulate every stage of the network through spatial feature
transforms~\cite{wang2018sft}, since the regime has to be known at every depth.
\paragraph{Network.} $G$ is a four-level U-Net. The two finest levels use convolutional NAFNet-style
blocks~\cite{chen2022nafnet}; at $1/4$ resolution channel attention~\cite{zamir2022restormer} mixes
information across the tile at linear cost; at $1/8$ resolution ($64\times64$ tokens per tile) spatial
self-attention provides tile-wide context and cross-attention reads the exemplar memory.
\todo{sizes; whether a larger model initialized from a pretrained super-resolution network~\cite{chen2023hat}
is reported (used to validate the formulation).}
\paragraph{Losses.} An $\ell_1$ loss on the full output dominates in focus. The f/22 target is slightly
misregistered and has its own noise, so pixel losses alone would favor the mean and treat transferred
texture as noise; in defocused regions we add an alignment-tolerant texture
loss~\cite{mechrez2018contextual} and, in a second stage, a light adversarial loss.
\todo{weights, schedule.}
\paragraph{Generative variant.} Where the memory abstains under strong defocus, a generative prior
can fill in. We adapt a one-step super-resolution model~\cite{wu2026vosr}, the only generative
upsampler in our study that leaves in-focus regions intact, with LoRA, conditioned on $x_T$, $d_T$ and
the exemplar tokens, inside the anchor lock of \cref{eq:output}, with the noise shared across tiles.
\todo{only if the feed-forward model leaves a visible gap in such regions; otherwise an ablation.}

\subsection{Training data and plug-in anchors}
\label{sec:training}
We train on DPDD at native resolution ($350$ scenes, \cref{sec:exp_data}). The upsampler sees $\mathcal{D}$
only through the anchor, so any deblurring method can be plugged in at test time. Most published
methods were trained on DPDD at exactly the anchor resolution, which makes their training-set anchors
unrealistically good; Bokehlicious~\cite{seizinger2025bokehlicious} was not, and provides our main
training anchors. They are complemented by anchors of other methods and by simulated anchors (the
downscaled target degraded by residual blur that grows with $d$, ringing and noise). Methods held out
of training test the plug-in claim.
